\let\textcircled=\pgftextcircled
\chapter*{Annexe}
\label{chap:intro}
\addstarredchapter{Annexe}
\begin{enumerate}
\item \textbf{Présentation du Fonds Routier}\\
Le Fonds routier a été créé par la loi $n\circ 96/07$  du 08 avril 
 1996 portant protection du patrimoine national et est régit par le décret $N°2005/239$ du 24 juin 2005qui en fixe les modalités de fonctionnement. C’est un établissement public administratif (EPA) dédié au financement de l’entretien routier au Cameroun. Le Fonds routier assure le financement, d’une part, des programmes de protection du patrimoine routier national, des programmes de prévention et de sécurité routières, d’entretien du réseau routier, et d’autre part, des opérations de réhabilitation et d’aménagement des routes. Etant un établissement public administratif de type particulier par rapport à ses organes de gestion, à la rémunération et aux avantages de son personnel, ainsi qu’aux règles de tenue de sa comptabilité, le Fonds routier est doté de la personnalité juridique et jouit de l’autonomie de gestion.
\vspace{1cm}
\item \textbf{Mission du Fonds routier}\\
Le Fonds routier exerce ses missions de financement de l’entretien routier et de
paiement des prestations réalisées à l’entreprise à travers deux guichets distincts à savoir, le
guichet «entretien» et le guichet «investissement». Le guichet« Entretien »a pour objectif exclusif d’assurer le financement et le paiement des prestations réalisées par l’entreprise et relatives à trois aspects liés à la sécurité. Quant au Le guichet« Investissement », il a pour objectif exclusif d’assurer le financement et le paiement des prestations réalisées au titre de
l’aménagement et à la réhabilitation des routes.
\vspace{1cm}
\item \textbf{Ressources du Fonds routier}\\
Les ressources du Fonds routier proviennent:\\
\textbf{Pour le guichet Entretien} :
\begin{itemize}
    \item  La redevance d’usage de la route (RUR);
 Le droit de péage routier, ou en cas de concession du péage, de la redevance de
concession;
 \item Les dotations budgétaires des Ministères destinées à alimenter la ligne d’urgence au titre des interventions d’urgence;
\item Les ressources provenant des produits financiers générés par le placement des
excédents de trésorerie éventuels;
\item Le produit de la taxe à l’essieu
\item Le produit de la taxe de transit;
\item Le produit des amendes
\end{itemize}


\textbf{Pour le guichet investissement };
\begin{itemize}
    \item Les dons, legs, subventions et aides diverses apportés par les partenaires
financiers du Cameroun;
\item Les dotations budgétaires des Ministères;
\end{itemize}
\end{enumerate}

